\section{Conclusion}
\label{sec:conclusion}

A corrupted-image difference between quantized detectors is not, by itself, a
robustness comparison. It combines any pre-existing clean-fidelity gap with the
additional change induced by corruption. We introduced a byte-matched,
image-paired, floor-guarded protocol that separates these quantities through the
direct interaction
$\Delta E=(\mathrm{FP8}-\mathrm{INT8})_c-
(\mathrm{FP8}-\mathrm{INT8})_{95}$ and interprets it only alongside matched-clean
and absolute corrupted AP.

The distinction changes the evidence. Although the raw corrupted gap favored
FP8 in 134 of 144 YOLO cells, clean adjustment changed the interpretation in 56
cells. FP8 preserved 1.40 more matched-clean AP on average across all 12 YOLO
blocks, whereas the balanced exploratory interaction was -0.02 AP points with a
95\% interval of $[-0.24,+0.15]$. Untouched VOC/KITTI final partitions yielded
-0.55 $[-0.78,-0.29]$, demonstrating a narrower but stronger split-independent
conditional result. Absolute accuracy fell in 282 of 288 corrupted arms, with
35 below 10 AP, so gap contraction cannot be equated with retained performance.
Architecture stress cases further showed that a negative interaction can arise
when an INT8 recipe already fails its matched-clean fidelity gate.

The supported conclusion is not a universal INT8--FP8 robustness ranking.
Clean fidelity, absolute corrupted accuracy, and the corruption-specific format
interaction are different deployment claims. Reporting them together, under an
explicit finite target and a fully specified executable treatment, produces a
more auditable and decision-relevant evaluation of quantized computer-vision
systems.
